% ============================================================================
% DHG508 Final Design Paper & Post-Mortem (Simplified Overleaf Template)
% Lingnan University — Critical Digital History
%
% OVERLEAF SETUP:
%   1. Upload dhg508-overleaf-template.zip to Overleaf (New Project -> Upload Project)
%   2. Click Menu (top left) -> Compiler -> select XeLaTeX
%   3. Click Recompile (Biber runs automatically for Chicago footnotes & bibliography)
% ============================================================================

\documentclass[12pt,a4paper]{article}

\usepackage{fontspec}
\newcommand{\MainLatinFont}{Times New Roman}
\newcommand{\FallbackLatinFont}{TeX Gyre Termes}
\newcommand{\MainCJKFont}{Noto Serif CJK SC}
\newcommand{\FallbackCJKFont}{FandolSong-Regular.otf}

\IfFontExistsTF{\MainLatinFont}
  {\setmainfont{\MainLatinFont}}
  {\setmainfont{\FallbackLatinFont}}

\usepackage{xeCJK}
\IfFontExistsTF{\MainCJKFont}
  {\setCJKmainfont{\MainCJKFont}}
  {\IfFontExistsTF{Songti SC}
    {\setCJKmainfont{Songti SC}}
    {\IfFontExistsTF{SimSun}
      {\setCJKmainfont{SimSun}}
      {\setCJKmainfont{\FallbackCJKFont}}}}

\usepackage[a4paper, margin=1in]{geometry}
\usepackage{setspace}
\usepackage{booktabs}
\usepackage{array}
\usepackage{graphicx}
\usepackage{csquotes}

% Chicago Manual of Style 18th ed., Notes-Bibliography (BibLaTeX + Biber)
\usepackage[notes, backend=biber, autolang=other, isbn=false]{biblatex-chicago}
\addbibresource{references.bib}

\renewcommand{\footnotesize}{\fontsize{10}{12}\selectfont}
\usepackage[hidelinks]{hyperref}

\title{\textbf{Replace With Your Project Title}\\[0.3em]
\large Design Rationale, Evidence Boundaries, and Post-Mortem}
\author{Your Name\\Lingnan University $\cdot$ DHG508 Critical Digital History\\
\small Live URL: \url{https://your-project.pages.dev} $\cdot$ Repo: \url{https://github.com/your-name/your-repo}}
\date{25 November 2026}

\begin{document}
\maketitle

\begin{abstract}
\noindent
Replace this paragraph (100--150 words) with your own summary: state the historical
or cultural question your project investigates, what a user does inside your
application or simulation, how documented facts (\emph{据}), historical inferences
(\emph{推断}), and design inventions (\emph{设}) are separated, and the main lesson
from building and testing it with AI agents.
\end{abstract}

\onehalfspacing

\section{Historical Question and Why This Medium Matters}

State your core historical or cultural question and your design question. Explain
why an interactive simulation, database application, or spatial reconstruction
teaches something that a static essay or film would not---for example, what a user
discovers only by choosing, testing constraints, or navigating the space.

\section{Sources, Limits, and Scholarly Conversation}

Describe the primary sources your project rests on (with exact editions and page
locators) and state their limits plainly (translation layers, OCR errors, missing
voices, or archival gaps).\autocite[150, 152]{tobin1986macro} Use \verb|\autocite[page]{key}|
after punctuation to generate Chicago footnotes automatically from \texttt{references.bib}.\autocite{conference1986}

Engage with \textbf{at least three scholarly works} relevant to your problem,
explaining what those scholars argue, where they agree or disagree, and how your
design choices respond to that conversation.\autocite[158--160]{gewirtz2017unlikely}
Connect each scholarly work directly to a concrete mechanism, character rule, or
data structure in your build.\autocite[195--199]{weber2021shock}

\section{System Architecture and Epistemic Boundaries}

Explain how your system works end to end: how data (\texttt{.csv}/\texttt{.json})
and rules are separated from UI or engine rendering, how live model calls (if any)
are proxied safely via a Cloudflare Worker or pre-computed offline, and what hard
temporal or epistemic cutoffs keep characters or outputs from drifting into
anachronism.

% Uncomment the figure block below once you upload schematic.png alongside main.tex:
% \begin{figure}[htbp]
%   \centering
%   \includegraphics[width=0.92\textwidth]{schematic.png}
%   \caption{Project architecture and interaction flow generated during ideation.}
%   \label{fig:schematic}
% \end{figure}

\section{Reproducible Evidence Table}

Table~\ref{tab:evidence} is generated directly from \texttt{evidence.csv} by
\texttt{make\_evidence\_table.py} so that your repository data and this paper share
a single source of truth.

\begin{table}[htbp]
  \centering
  \footnotesize
  \singlespacing
  \renewcommand{\arraystretch}{1.25}
  % Auto-generated from evidence.csv by make_evidence_table.py — edit the CSV, not this file.
\begin{tabular}{p{0.27\textwidth}p{0.20\textwidth}p{0.13\textwidth}p{0.08\textwidth}p{0.22\textwidth}}
  \toprule
  \textbf{Claim / Element} & \textbf{Source / File} & \textbf{Locator} & \textbf{Level} & \textbf{Boundary / Limit} \\
  \midrule
  Tobin: demand management outcomes are never certain & Bashan Conference Volume (1986) & p. 150 & 据 & Printed Chinese translation; does not reject reform \\
  Indirect monetary levers require responsive banks and statistics & Bashan Conference Volume (1986) & pp. 152, 157-158 & 据 & On-the-spot Chinese reply is not recorded in the volume \\
  Debate tension pairs Western preconditions with Q4 enterprise credit pressure & Weber (2021); Xiang Huaicheng memoir & p. 195; P6, P11 & 推断 & Historian's synthesis across two separate texts \\
  Bow-deck encounter between James Talbot and Cheng Xiaonong & game/people/expert-1.json & zone: bashan-anchor & 设 & Ship deck is documented; this specific dialogue is invented \\
  \bottomrule
\end{tabular}

  \caption{Core claims and design elements with locators, evidence tiers
  (\emph{据} documented, \emph{推断} inferred, \emph{设} invented), and boundaries.}
  \label{tab:evidence}
\end{table}

\section{Automated Verification and External Playtesting}

Summarize what your automated test suite (\texttt{Vitest}/\texttt{pytest},
\texttt{Playwright} E2E, and any log audit) checks across both software behaviour
and historical constraints. Report what at least \textbf{two external users}
experienced on your deployed URL, what broke or confused them, and which permanent
regression tests you added so those issues cannot return.

\section{Post-Mortem and Statement of AI Use}

Give a concrete, candid account grounded in your actual research journal and Git
history:
\begin{enumerate}
  \item \textbf{What worked:} Which steps (search triage, schema refactoring,
  schematic generation, Browser MCP debugging, test generation) genuinely helped?
  \item \textbf{What failed or drifted:} Give specific examples of OCR mistakes,
  goal drift, hallucinated citations, broken builds, or deployment snags, and how
  you caught and resolved them.
  \item \textbf{What changed:} How your current build differs from your initial
  idea and why.
  \item \textbf{What this work does \emph{not} claim:} Honest historical and
  technical limits that remain open.
  \item \textbf{AI disclosure:} Which models, agent harnesses, and MCP tools were
  used at which stages, and how you verified every source and claim yourself.
\end{enumerate}

\singlespacing
\printbibliography[title={Bibliography}]

\end{document}
